%%%%%%%%%%%%%%%%%%%%%%%%%%%%%%%%%%%%%%%%%%%%%%%%%%%%%%%%%%%%%%%%%%%%%%%%%%%%%%%%
% ACF-DOC-039.05
% Wildfire Prediction System
%%%%%%%%%%%%%%%%%%%%%%%%%%%%%%%%%%%%%%%%%%%%%%%%%%%%%%%%%%%%%%%%%%%%%%%%%%%%%%%%

\section{Wildfire Prediction System}

\subsection{Executive Summary}

The ACF Wildfire Prediction System is a global environmental intelligence
module dedicated to wildfire detection, prediction, risk assessment and
operational response.

Wildfires are complex Earth system events resulting from interactions
between:

\begin{itemize}

\item Atmospheric conditions

\item Vegetation state

\item Soil moisture

\item Land surface properties

\item Human activities

\item Climate variability

\end{itemize}


The ACF Wildfire System integrates:

\begin{itemize}

\item Numerical Weather Prediction

\item Fire danger indices

\item Satellite observations

\item Land surface models

\item Artificial intelligence

\item Earth System Digital Twin

\item Emergency management systems

\end{itemize}


The objectives are:

\begin{itemize}

\item Early wildfire detection

\item Fire danger forecasting

\item Fire propagation simulation

\item Smoke impact analysis

\item Emergency decision support

\end{itemize}


%%%%%%%%%%%%%%%%%%%%%%%%%%%%%%%%%%%%%%%%%%%%%%%%%%%%%%%%%%%%%%%%%%%%%%%%%%%%%%%%

\subsection{Introduction}

Wildfires represent one of the most important environmental hazards
affecting ecosystems, populations and atmospheric composition.

A wildfire requires three fundamental elements:

\[
Fire =
Fuel
+
Oxygen
+
Heat
\]


The ACF system analyzes the complete fire environment:

\[
Weather
+
Vegetation
+
Terrain
+
HumanFactors
\rightarrow
FireRisk
\]


The wildfire prediction chain is:

\[
Observation
\rightarrow
RiskAnalysis
\rightarrow
Detection
\rightarrow
PropagationModel
\rightarrow
Warning
\]


%%%%%%%%%%%%%%%%%%%%%%%%%%%%%%%%%%%%%%%%%%%%%%%%%%%%%%%%%%%%%%%%%%%%%%%%%%%%%%%%

\subsection{Wildfire Earth System Concept}

The ACF wildfire module connects multiple Earth components:

\begin{itemize}

\item Atmosphere

\item Biosphere

\item Land surface

\item Carbon cycle

\item Atmospheric chemistry

\item Human exposure

\end{itemize}


The system represents wildfire as a coupled
atmosphere--land--chemistry process.


%%%%%%%%%%%%%%%%%%%%%%%%%%%%%%%%%%%%%%%%%%%%%%%%%%%%%%%%%%%%%%%%%%%%%%%%%%%%%%%%

\subsection{Fire Weather Analysis}

Weather conditions strongly influence wildfire probability.

The ACF Fire Weather Engine monitors:

\begin{itemize}

\item Temperature

\item Relative humidity

\item Wind speed

\item Wind direction

\item Precipitation deficit

\item Atmospheric stability

\end{itemize}


High wildfire danger occurs when:

\[
FireDanger
=
f(T,RH,W,P)
\]


where:

\begin{itemize}

\item $T$ = temperature

\item $RH$ = relative humidity

\item $W$ = wind

\item $P$ = precipitation conditions

\end{itemize}


%%%%%%%%%%%%%%%%%%%%%%%%%%%%%%%%%%%%%%%%%%%%%%%%%%%%%%%%%%%%%%%%%%%%%%%%%%%%%%%%
%%%%%%%%%%%%%%%%%%%%%%%%%%%%%%%%%%%%%%%%%%%%%%%%%%%%%%%%%%%%%%%%%%%%%%%%%%%%%%%%
\subsection{Fire Danger Index}

The ACF Fire Danger Engine calculates wildfire probability using
meteorological and environmental variables.

The system integrates classical fire danger indices with artificial
intelligence corrections.


The main indicators are:

\begin{itemize}

\item Canadian Fire Weather Index (FWI)

\item Keetch-Byram Drought Index (KBDI)

\item Burning Index (BI)

\item Energy Release Component (ERC)

\item Fire Potential Index (FPI)

\end{itemize}


%%%%%%%%%%%%%%%%%%%%%%%%%%%%%%%%%%%%%%%%%%%%%%%%%%%%%%%%%%%%%%%%%%%%%%%%%%%%%%%%

\subsubsection{Canadian Fire Weather Index}

The Fire Weather Index evaluates the potential intensity of wildfire
development.

The general formulation is:

\[
FWI=f(M_f,T,RH,W)
\]


where:

\begin{itemize}

\item $M_f$ = fuel moisture

\item $T$ = temperature

\item $RH$ = relative humidity

\item $W$ = wind speed

\end{itemize}


The ACF implementation calculates:

\begin{itemize}

\item Fine fuel moisture

\item Duff moisture

\item Drought conditions

\item Fire spread potential

\end{itemize}


%%%%%%%%%%%%%%%%%%%%%%%%%%%%%%%%%%%%%%%%%%%%%%%%%%%%%%%%%%%%%%%%%%%%%%%%%%%%%%%%

\subsection{Fuel Moisture Monitoring}

Fuel moisture is a key parameter controlling wildfire ignition.

The ACF system monitors:

\begin{itemize}

\item Dead vegetation moisture

\item Live vegetation moisture

\item Soil moisture

\item Leaf water content

\item Biomass availability

\end{itemize}


Fuel dryness increases when:

\[
FuelDryness
=
Temperature
-
Humidity
-
Precipitation
\]


The model receives data from:

\begin{itemize}

\item Satellite observations

\item Land surface models

\item Meteorological forecasts

\item Vegetation models

\end{itemize}


%%%%%%%%%%%%%%%%%%%%%%%%%%%%%%%%%%%%%%%%%%%%%%%%%%%%%%%%%%%%%%%%%%%%%%%%%%%%%%%%

\subsection{Vegetation and Biomass Analysis}

Vegetation represents the available fuel for wildfire propagation.

The ACF biosphere module evaluates:

\begin{itemize}

\item Vegetation type

\item Biomass density

\item Leaf area index

\item Vegetation health

\item Carbon content

\end{itemize}


Important parameters:

\[
FuelLoad=f(Biomass,Vegetation,Moisture)
\]


The system generates:

\begin{itemize}

\item Fuel availability maps

\item Vegetation stress maps

\item Ignition probability maps

\end{itemize}


%%%%%%%%%%%%%%%%%%%%%%%%%%%%%%%%%%%%%%%%%%%%%%%%%%%%%%%%%%%%%%%%%%%%%%%%%%%%%%%%

\subsection{Fire Ignition Probability}

Wildfire ignition can result from natural and human causes.

The ACF ignition model considers:

\begin{itemize}

\item Lightning activity

\item Temperature extremes

\item Dry vegetation

\item Human activity

\item Wind conditions

\end{itemize}


The ignition probability is:

\[
P_{fire}
=
f(Weather,Fuel,Ignition)
\]


%%%%%%%%%%%%%%%%%%%%%%%%%%%%%%%%%%%%%%%%%%%%%%%%%%%%%%%%%%%%%%%%%%%%%%%%%%%%%%%%

\subsection{Fire Propagation Modeling}

Once ignition occurs, fire evolution depends on environmental
conditions.

The ACF Fire Dynamics Model simulates:

\begin{itemize}

\item Fire front movement

\item Rate of spread

\item Fire intensity

\item Burned area expansion

\end{itemize}


The fire spread velocity is represented by:

\[
V_f=f(W,S,F,M)
\]


where:

\begin{itemize}

\item $W$ = wind

\item $S$ = slope

\item $F$ = fuel properties

\item $M$ = moisture

\end{itemize}


%%%%%%%%%%%%%%%%%%%%%%%%%%%%%%%%%%%%%%%%%%%%%%%%%%%%%%%%%%%%%%%%%%%%%%%%%%%%%%%%

\subsection{Terrain Influence}

Topography strongly affects wildfire behavior.

The ACF terrain model includes:

\begin{itemize}

\item Elevation

\item Slope

\item Aspect

\item Terrain channels

\item Mountain effects

\end{itemize}


Fire usually propagates faster uphill due to:

\begin{itemize}

\item Increased heating

\item Better fuel exposure

\item Wind acceleration

\end{itemize}


%%%%%%%%%%%%%%%%%%%%%%%%%%%%%%%%%%%%%%%%%%%%%%%%%%%%%%%%%%%%%%%%%%%%%%%%%%%%%%%%
%%%%%%%%%%%%%%%%%%%%%%%%%%%%%%%%%%%%%%%%%%%%%%%%%%%%%%%%%%%%%%%%%%%%%%%%%%%%%%%%
\subsection{Satellite Wildfire Detection}

Satellite observation provides global and continuous wildfire
monitoring capability.

The ACF Earth Observation System integrates:

\begin{itemize}

\item Thermal infrared observations

\item Optical satellite imagery

\item Vegetation monitoring

\item Burned area products

\item Aerosol observations

\end{itemize}


The main satellite sources include:

\begin{itemize}

\item MODIS fire products

\item VIIRS active fire detection

\item Sentinel optical observations

\item Geostationary satellite monitoring

\end{itemize}


Satellite observations provide:

\begin{itemize}

\item Active fire location

\item Fire intensity estimation

\item Burned area mapping

\item Fire evolution monitoring

\end{itemize}


%%%%%%%%%%%%%%%%%%%%%%%%%%%%%%%%%%%%%%%%%%%%%%%%%%%%%%%%%%%%%%%%%%%%%%%%%%%%%%%%

\subsection{Thermal Anomaly Detection}

Active fires generate strong thermal signatures.

The ACF detection system analyzes:

\begin{itemize}

\item Land surface temperature anomalies

\item Infrared radiation

\item Thermal contrast

\item Night-time fire signals

\end{itemize}


The thermal anomaly is:

\[
A_T=
T_{fire}
-
T_{background}
\]


where:

\begin{itemize}

\item $T_{fire}$ = fire temperature

\item $T_{background}$ = surrounding temperature

\end{itemize}


%%%%%%%%%%%%%%%%%%%%%%%%%%%%%%%%%%%%%%%%%%%%%%%%%%%%%%%%%%%%%%%%%%%%%%%%%%%%%%%%

\subsection{Smoke and Atmospheric Chemistry Impact}

Wildfires modify atmospheric composition.

The ACF Atmospheric Chemistry Module simulates:

\begin{itemize}

\item Smoke transport

\item Aerosol emissions

\item Carbon monoxide

\item Carbon dioxide

\item Particulate matter

\end{itemize}


Fire emissions are represented as:

\[
Emission=
BurnedArea
\times
FuelLoad
\times
EmissionFactor
\]


The system predicts:

\begin{itemize}

\item Air quality degradation

\item Visibility reduction

\item Population exposure

\item Long-range smoke transport

\end{itemize}


%%%%%%%%%%%%%%%%%%%%%%%%%%%%%%%%%%%%%%%%%%%%%%%%%%%%%%%%%%%%%%%%%%%%%%%%%%%%%%%%

\subsection{Artificial Intelligence Fire Detection Engine}

The ACF AI Fire Engine improves wildfire detection and forecasting.

The hybrid architecture combines:

\[
Prediction=
PhysicsModel
+
SatelliteData
+
AI
\]


The AI system performs:

\begin{itemize}

\item Satellite image classification

\item Fire hotspot detection

\item Smoke recognition

\item Fire spread prediction

\item Damage estimation

\end{itemize}


%%%%%%%%%%%%%%%%%%%%%%%%%%%%%%%%%%%%%%%%%%%%%%%%%%%%%%%%%%%%%%%%%%%%%%%%%%%%%%%%

\subsection{Deep Learning Fire Models}

The wildfire AI framework includes:

\begin{itemize}

\item Convolutional Neural Networks for satellite images

\item Vision Transformers for Earth observation

\item Recurrent networks for fire evolution

\item Graph Neural Networks for landscape connectivity

\end{itemize}


Input variables:

\begin{itemize}

\item Satellite imagery

\item Weather forecasts

\item Wind fields

\item Vegetation maps

\item Terrain information

\end{itemize}


AI outputs:

\begin{itemize}

\item Fire probability maps

\item Fire perimeter prediction

\item Spread uncertainty

\item Impact estimation

\end{itemize}


%%%%%%%%%%%%%%%%%%%%%%%%%%%%%%%%%%%%%%%%%%%%%%%%%%%%%%%%%%%%%%%%%%%%%%%%%%%%%%%%

\subsection{Earth System Digital Twin Integration}

The ACF Wildfire System is integrated into the Earth System Digital Twin.

The Digital Twin provides:

\begin{itemize}

\item Atmospheric forecast

\item Land surface state

\item Vegetation condition

\item Carbon cycle information

\item Climate scenarios

\end{itemize}


The coupling enables:

\begin{itemize}

\item Real-time fire simulation

\item Future wildfire scenarios

\item Climate change impact analysis

\item Ecosystem recovery monitoring

\end{itemize}


%%%%%%%%%%%%%%%%%%%%%%%%%%%%%%%%%%%%%%%%%%%%%%%%%%%%%%%%%%%%%%%%%%%%%%%%%%%%%%%%

